\documentclass[../main.tex]{subfiles}

\begin{document}

\chapter{Discussion}\label{ch:discussion}

This chapter discusses the limitations encountered during this study and explores potential 
future implications for research and tool development in the area of application layer 
monitoring tool for IdPs. 

\section{Interpretation of the Results}\label{sec:interpretation}

The results of this study can be read along the three research questions posed in 
Chapter~\ref{ch:introduction}. 
The first question asked which security relevant data can be collected at the application layer 
of a Shibboleth IdP and through which mechanisms. 
Three sources proved usable, namely the audit log, the configuration and metadata of the SPs, 
and the runtime state accessible through an interceptor. 
The value of this combination does not primarily result from the number of sources but from the 
positions at which they observe the authentication process. 
The audit log records the outcome of a completed transaction, the configuration states what was 
intended for the partner involved, and the interceptor observes the incoming message while it is 
being processed. 
Each source therefore holds information that the others cannot provide, and the complementarity 
follows from their position in the processing sequence rather than from a deliberate division of 
responsibilities. 

The second question asked how these data can be correlated in order to detect configuration 
weaknesses and anomalies. 
The correlation operates reliably within a limit that the evaluation was able to quantify, namely 
approximately six logins per second for a given combination of subject, address and SP, 
beyond which the heuristic strategies can no longer resolve the order of events. 
The thirteen detection rules behave as intended for the scenarios constructed for them, and all 
are verified against real persistence rather than against mocked repositories alone. 
This establishes the fidelity of the implementation, though not the detection quality in 
operation, a distinction developed in Section~\ref{subsec:validity-detection}.

The third question asked what overhead the resulting system imposes. 
In the asynchronous delivery mode no overhead was measurable, with the work performed inside the 
interceptor amounting to approximately 0,3\% of the login duration. 
The synchronous mode tends to be slower by a margin in the low tens of milliseconds. 

Beyond the individual research questions, one observation emerges across all three evaluations. 
In each of them the measurement uncovered a defect that operation itself had not revealed. 
A missing version field caused every behaviour profile update to be discarded silently, a 
pipeline defect would have reported every encrypted login as unencrypted, and a misconfigured 
audit format left one correlation strategy permanently inactive. 
In none of these cases did the system fail visibly. 
It continued to run, alerts continued to be produced, and the dashboard continued to be 
populated. 
The defects became apparent only when the behaviour was compared against an expectation defined 
in advance.

This observation is worth stating separately because it bears on the premise of the work itself. 
The argument for observing an IdP from within rests on the assumption that a system which appears 
to function is not thereby known to function. 
The three defects are instances of exactly that, encountered in the monitoring system rather than 
in the monitored one, which does not weaken the argument but illustrates its generality.

\section{Limitation of the Study}\label{sec:limitation}

This section discusses the structural limitations of the proposed monitoring system. 
Rather than merely listing limitations, it analyses their underlying causes and shows how they result 
from the position of the monitoring system within the IdP.
As briefly discussed in Section~\ref{subsec:detection-attack-space}, the system addresses two of the 
six SAML cross-protocol issues identified in~\cite{FIMSurvey}. 
This limited coverage is not primarily caused by missing implementations, but by the observation position 
of the monitoring system, which operates from within the IdP. 
The remaining issues can be categorised into three groups: a different system level, a different party, 
and indirect observability. 

The first group consists of issues that primarily affect a different system layer. 
DNS poisoning, for example, affects the network layer below the application layer. 
The IdP processes the request it receives regardless of how the client resolved the corresponding address. 
The same applies to DDoS and UI Redressing, which occur outside the application-level observation point 
of the IdP.

The second group consists of issues that affect a different party within the federation. 
\acrfull{csrf}, for example, describes an implementation vulnerability at the SP. 
The relevant interaction occurs between the browser and the SP and therefore does not pass through the IdP. 
The same applies to XSS and Session Swapping.

The third group contains issues that can only be observed indirectly from the IdP. 
A \acrfull{mitm} attack is an example of this category. 
The attack occurs on the communication path after the message has left the IdP. 
The monitoring system can therefore observe prerequisites that make such an attack possible, but not the 
attack itself.
For example, if an assertion is transmitted without encryption, an attacker could potentially read or modify 
the message in transit. 
The monitoring system can detect that the assertion was transmitted without encryption, but it cannot 
determine from the IdP whether the message was actually intercepted or modified.

Two cases are considered separately because of their particular relevance.
The first is the \textit{Bogus Merchant} scenario. 
This is the only cross-protocol issue identified in the survey that directly concerns Shibboleth, but it is 
not covered by the monitoring system. 
The reason is that the malicious SP is a regular federation participant and is registered in the federation 
metadata as a legitimate SP. 
From the perspective of the IdP, the SP therefore does not necessarily exhibit an observable property that 
distinguishes it from other legitimate participants. 
Monitoring changes to SP configuration and metadata could provide an additional source of information for 
detecting suspicious changes, but this alone cannot establish that a formally registered SP is malicious.

The second case is more fundamental. 
The \textit{Malicious Provider} scenario describes an IdP that abuses its position within the federation. 
Since the monitoring system is directly integrated into the IdP, the observer and the potential attacker are 
located within the same system. 
An entity that fully controls the IdP can therefore also control the data collection layer of the monitoring 
system. 
This limitation follows directly from the observation point and cannot be overcome by adding further 
detection rules within the same trust domain. 

\section{Implication for Future Research}\label{sec:implication}

Was aus den Grenzen folgt. Die Bogus-Merchant-Zeile über Metadaten-Änderungserkennung teilweise adressierbar. Der outboundMessageFlow für die Felder, die heute nur im Audit-Log stehen. SP-seitige Erkennung als Ergänzung, mit Verweis auf Martín und Beltrán. OIDC analytisch: Gilt die Drei-Quellen-Architektur dort auch, und welche Regeln hätten eine Entsprechung.

Und was ohne neue Datengrundlage nicht geht: Erkennungsgüte im Betrieb braucht einen gelabelten Produktivdatensatz, den niemand hat.

\end{document}